\chapter{Robot Platform and Integration}
\label{cap:platform}

This chapter describes the hardware and software components of the robot platform used for field validation. It covers the camera configuration, computing hardware, software architecture using ROS, and the deployment of the trained segmentation model for real-time inference.

\section{Robot hardware and sensors}
\label{sec:hardware}

\subsection{Robot platform}
\label{subsec:robot_platform}

\todobox{Specifiche della piattaforma robotica mobile}{%
Descrivere il modello del robot mobile impiegato per le prove (es. piattaforma UGV a guida differenziale o 4WD):
\begin{itemize}
  \item Dimensioni fisiche d'ingombro (lunghezza, larghezza di carreggiata per il passaggio tra i filari, altezza da terra del telaio).
  \item Tipologia di trazione, attuatori elettrici, ruote motrici e ruote pivottanti di supporto.
  \item Peso totale, capacità di carico utile (payload) e velocità operativa nominale di avanzamento.
\end{itemize}
}

\subsection{Camera configuration}
\label{subsec:camera}

\placeholderfigure{fig:robot_camera_setup}%
{Physical mobile robot platform and front camera mounting setup}%
{Fotografia della piattaforma mobile con indicazione della telecamera RGB frontale, dell'altezza di montaggio $h$ dal terreno, dell'angolo di inclinazione (pitch) verso il suolo e dell'alloggiamento del calcolatore di bordo.}%
{5.5cm}

\todobox{Specifiche e geometria di montaggio della telecamera}{%
Dettagliare i parametri ottici e geometrici del sensore visivo:
\begin{itemize}
  \item Modello della telecamera RGB (es. Intel RealSense D435 / telecamera USB industriale).
  \item Risoluzione di acquisizione attiva ($1920\times 1080$ o $640\times 480$) e framerate di streaming nativo (30 fps).
  \item Posizione di montaggio: altezza dal suolo $h$, angolo di tilt verso il basso $\alpha$, e campo visivo orizzontale/verticale (HFOV, VFOV).
  \item Lente e lunghezza focale, formato di compressione dello stream in ingresso (JPEG compresso vs raw).
\end{itemize}
}

\subsection{Computing platform}
\label{subsec:computing}

\todobox{Unit\`a di calcolo di bordo NVIDIA Jetson e alimentazione}{%
Specificare l'hardware per l'elaborazione a bordo:
\begin{itemize}
  \item Modello dell'unit\`a embedded (es. NVIDIA Jetson Xavier NX / Orin), architettura CPU (core ARM Cortex) e GPU (Tensor Core).
  \item Memoria RAM unificata (GB) e supporto per l'accelerazione CUDA.
  \item Sistema operativo (Ubuntu con pacchetto NVIDIA JetPack) e ambiente di esecuzione runtime.
  \item Sistema di alimentazione elettrica: pacco batterie dedicato, convertitore di tensione DC-DC, autonomia oraria e consumo medio in watt durante l'inferenza.
\end{itemize}
}

\section{Software architecture}
\label{sec:software}

\subsection{ROS integration}
\label{subsec:ros}

The software pipeline is implemented using the Robot Operating System (ROS). ROS provides a modular architecture where perception, geometric extraction, and control components communicate through topics and services.

\placeholderfigure{fig:ros2_architecture}%
{ROS 2 software architecture node graph and topic communication flow}%
{Diagramma a blocchi dei nodi ROS 2 (\texttt{camera\_node}, \texttt{controller\_node}, \texttt{motor\_driver}), con evidenziati i topic (\texttt{/camera/image\_raw}, \texttt{/path\_mask}, \texttt{/cmd\_vel}) e i relativi messaggi.}%
{5cm}

\todobox{Architettura dei nodi ROS 2 e politiche QoS}{%
Descrivere l'organizzazione software nel package ROS 2 \texttt{seg\_controller\_onnx}:
\begin{itemize}
  \item Grafo dei nodi attivi e rispettivi ruoli nel ciclo di percezione-azione.
  \item Topic scambiati: ricezione stream telecamera compresso, pubblicazione della maschera del sentiero, coordinate della linea di navigazione e comandi motore \texttt{/cmd\_vel}.
  \item Politiche di Quality of Service (QoS): adozione di code a profondit\`a unitaria (\texttt{depth=1}) con policy \texttt{BEST\_EFFORT} per evitare l'accumulo di latenza su frame arretrati.
\end{itemize}
}

\subsubsection{Perception node}
\label{subsubsec:perception_node}

The perception node receives camera images, runs the segmentation model, and publishes the path mask. Input images are preprocessed according to the training configuration before inference.

\todobox{Dettagli implementativi del nodo di percezione}{%
Descrivere le operazioni eseguite alla ricezione di ogni frame:
\begin{itemize}
  \item Decodifica dell'immagine compressa JPEG tramite \texttt{cv2.imdecode}.
  \item Ridimensionamento e normalizzazione del tensore d'ingresso (media e deviazione standard di ImageNet).
  \item Invocazione dell'inferenza tramite la sessione ONNX Runtime con \texttt{CUDAExecutionProvider}.
  \item Estrazione della maschera binaria di navigazione tramite operatore \texttt{argmax}.
\end{itemize}
}

\subsubsection{Geometric extraction node}
\label{subsubsec:geometric_node}

The geometric extraction node subscribes to the segmentation mask, applies Method~3 or Method~4 from Chapter~\ref{cap:geometric_reference}, and publishes the navigation reference as a line or position message.

\todobox{Integrazione dell'estratore geometrico e tracking temporale}{%
Dettagliare l'elaborazione geometrica:
\begin{itemize}
  \item Suddivisione della maschera in strisce orizzontali e selezione dei candidati.
  \item Applicazione del filtro particellare (\texttt{pf3.py} o \texttt{pf4.py}) per tracciare la linea centrale nel tempo, riducendo le oscillazioni introdotte da buchi nella vegetazione.
  \item Gestione dei frame non validi (reiezione in caso di confidenza o supporto insufficiente).
\end{itemize}
}

\subsubsection{Control node}
\label{subsubsec:control_node}

The control node converts the navigation reference into velocity commands using proportional control based on lateral displacement and heading error.

\todobox{Interfaccia di controllo e generazione comandi /cmd\_vel}{%
Specificare la logica di calcolo del comando:
\begin{itemize}
  \item Calcolo dell'offset laterale $e_y$ rispetto al centro del frame visivo e dell'inclinazione angolare $e_\theta$.
  \item Legge di controllo proporzionale per ricavare la velocit\`a lineare $v$ e la velocit\`a angolare $\omega$.
  \item Costruzione e pubblicazione del messaggio \texttt{geometry\_msgs/msg/Twist} sul topic \texttt{/cmd\_vel}.
\end{itemize}
}

\subsection{Perception pipeline implementation}
\label{subsec:perception_pipeline}

\todobox{Struttura a callback sincrona e carico delle immagini di debug}{%
Descrivere il funzionamento del nodo unificato in un unico thread:
\begin{itemize}
  \item Sequenzialit\`a del ciclo: decodifica JPEG $\to$ inferenza ONNX $\to$ particle filter $\to$ generazione comandi $\to$ compressione e pubblicazione di 3 immagini di diagnostica visiva.
  \item Analisi dell'impatto prestazionale: la compressione JPEG delle 3 immagini di debug rappresenta una frazione significativa del tempo totale di callback.
\end{itemize}
}

\subsection{Control interface}
\label{subsec:control_interface}

\todobox{Interfaccia con l'hardware di trazione del robot}{%
Descrivere come i messaggi \texttt{/cmd\_vel} vengono convertiti dai driver dei motori in comandi di velocit\`a angolare per le singole ruote motrici (tramite cinematica differenziale diretta e inversa).
}

\section{Deployment configuration}
\label{sec:deployment}

\subsection{Model export to ONNX}
\label{subsec:onnx_export}

The trained LR-ASPP model is exported to ONNX format for deployment. ONNX Runtime provides optimized inference across different hardware platforms. The export procedure preserves model weights and input/output specifications from the training checkpoint.

\todobox{Procedura di esportazione ONNX}{%
Documentare i comandi e le opzioni dello script di esportazione da PyTorch (.pt/.pth) a ONNX (.onnx):
\begin{itemize}
  \item Funzione \texttt{torch.onnx.export} con parametri di input (batch size fisso a 1, risoluzione d'ingresso $H \times W$, Opset 13/14).
  \item Verifica di equivalenza numerica tra i logit prodotti da PyTorch e quelli generati dal modello esportato in ONNX.
\end{itemize}
}

\subsection{Real-time inference}
\label{subsec:real_time_inference}

The deployed model processes camera frames at the rate required for closed-loop control. Input preprocessing matches the training transformations. The segmentation mask is decoded and passed to the geometric extraction stage.

\todobox{Configurazione del runtime ONNX}{%
Specificare i dettagli dell'ambiente di esecuzione:
\begin{itemize}
  \item Configurazione dell'allocazione di memoria e provider attivo (\texttt{CUDAExecutionProvider}).
  \item Tipologia di quantizzazione/precisione (es. FP32 standard vs FP16) e consumo di memoria VRAM sulla GPU Jetson.
\end{itemize}
}

\subsection{Timing and latency}
\label{subsec:timing}

End-to-end latency includes image acquisition, preprocessing, segmentation inference, geometric extraction, and control computation. The total pipeline latency must remain below the control update period to maintain stable navigation.

\placeholderfigure{fig:jetson_timing_breakdown}%
{Execution time breakdown of the synchronous perception-control cycle on NVIDIA Jetson}%
{Grafico della scomposizione del tempo di ciclo medio su Jetson: decodifica JPEG (10--15 ms), inferenza ONNX su GPU (10--15 ms), filtro geometrico/particellare (5--10 ms), codifica delle 3 immagini JPEG di debug (15--25 ms). Tempo totale per ciclo: 42--76 ms (~20 Hz).}%
{5cm}

\todobox{Analisi delle latenze misurate e frequenza a 20 Hz (dati reali)}{%
Riportare i risultati empirici ricavati dai test sul Jetson:
\begin{itemize}
  \item Tempo totale di esecuzione compreso tra 42 e 76 ms per ciclo (media $\approx 50$~ms), corrispondente a una frequenza di controllo di circa 20 Hz a fronte di una frequenza di cattura della telecamera di 30 fps.
  \item Analisi del frame dropping indotto dalla callback sincrona con \texttt{depth=1} e \texttt{BEST\_EFFORT} (viene scartato circa un frame ogni tre per mantenere il controllo sul dato pi\`u recente).
  \item Strategie per raggiungere i 30 Hz: (1) pubblicare le immagini di debug a frequenza ridotta o in un thread dedicato; (2) ottimizzare il modello tramite NVIDIA TensorRT; (3) disaccoppiare decodifica e inferenza con architettura multi-thread.
\end{itemize}
}
